\section{Introduction}\label{sec:intro}

Stars in dense relativistic clusters and collisionless dark matter around black holes move freely through collective gravitational fields that vary strongly in space and time \cite{1985ApJ...298...34S,2008gady.book.....B,Andreasson:2011ng,Angulo:2021kes,Bamber:2025gxj}. Two-body gravitational encounters act over timescales far longer than a dynamical crossing time, much as binary collisions remain negligible in a collisionless plasma \cite{1991ppcs.book.....B,2008gady.book.....B}. Multi-stream trajectories cross in configuration space and build anisotropic stresses, whereas single-fluid hydrodynamics assigns one bulk velocity to each spatial point \cite{Andreasson:2011ng,East:2019chx}. Strong-field evolution therefore requires the Einstein-Vlasov system, a phase-space kinetic description coupled to the full Einstein equations where matter follows geodesics of the curvature its stress-energy generates \cite{Andreasson:2011ng,Sarbach:2013fya,Acuna-Cardenas:2021nkj}. Early numerical studies with spherical codes mapped equilibrium sequences, radial stability, and black hole formation in relativistic clusters \cite{1968ApJ...154..251I,1980ApJ...238.1101I,1985ApJ...298...34S,1985ApJ...298...58S} alongside critical gravitational collapse \cite{Rein:1998uf,Olabarrieta:2001wy,Akbarian:2014gna}. Subsequent axisymmetric and three-dimensional simulations extended kinetic gravity to rotating clusters, nonspherical collapse, and cosmic structure \cite{Abrahams:1994ge,Shibata:1999va,Shibata:1999wi,East:2019bwu,East:2019chx,Ames:2020ssl}, where asymmetric instabilities and mergers demand evolution in three spatial dimensions.

The numerical-relativity particle-in-cell (NRPIC) method of Ref.~\cite{Wu:2026nrpic} brings algorithms from relativistic plasma simulation \cite{1988csup.book.....H,1991ppcs.book.....B,Spitkovsky:2005vsy,Cerutti:2013mma,EntityDevelopmentTeam:2025tlp} to three-dimensional Einstein-Vlasov evolution. Macroparticles---finite-mass packets sampling six-dimensional phase space---advance along timelike geodesics through a time-reversible trajectory pusher in the reconstructed metric \cite{boris1970relativistic,Bacchini:2020cze,Olivares:2021cfa,Boyeneni:2025tsx,Wu:2026nrpic}. Conservative cloud-in-cell deposition assigns each macroparticle's energy, momentum, and stress to nearby mesh cells across refinement boundaries before grid filtering smooths discrete sampling noise \cite{1988csup.book.....H,1991ppcs.book.....B,Myers:2021gfn,Wu:2026nrpic}. A constraint-damped solver in moving-puncture gauge, the standard way to evolve Einstein's equations through black hole formation on a fixed grid, advances the spacetime on an adaptive mesh \cite{Bona:2003fj,Campanelli:2005dd,Baker:2005vv,Hilditch:2012fp,Wu:2026nrpic}. Reference~\cite{Wu:2026nrpic} demonstrated this coupled cycle on gravitational collapse and relativistic cluster benchmarks using millions of macroparticles on large machines.

Two questions decide how broadly this method can be adopted. One question asks whether an independent builder without access to the original code can reconstruct the NRPIC method quantitatively from its published description alone. A numerical method is usable by others only when its written equations and algorithmic rules suffice for an independent implementation to recover its benchmarks. The second question asks which results survive at ten thousand macroparticles, where a run takes minutes to a few hours on one graphics card, and what must change for that. A small computational entry cost decides who can use the method for routine physical exploration, rapid parameter surveys, and independent verification.

Language-model agents under human direction built an independent implementation from the published text alone inside the public AthenaK code \cite{2026ApJS..283...27S,Zhu:2024utz}, with portable parallel execution through Kokkos \cite{2022ITPDS..33..805T}. Every algorithmic stage and coupled test was compared with exact analytical solutions or with the published numbers of Ref.~\cite{Wu:2026nrpic}. Candidate explanations of unexpected results were tested against predictions recorded before control runs designed so that they could fail, as described in Sec.~\ref{sec:assembly}.

The published description suffices to rebuild the method and reproduce the single-particle errors of Ref.~\cite{Wu:2026nrpic} once one unstated choice is identified (Sec.~\ref{sec:testparticles}). That choice is the tetrad gauge, the local rotation freedom of the orthonormal frame in the spacetime Boris pusher, where distinct frame constructions shift the discrete error (Secs.~\ref{sec:pushers} and \ref{sec:testparticles}). Self-gravitating evolutions at ten thousand macroparticles require two algorithmic changes (Sec.~\ref{sec:selfgravity}). Surface masses on a cubic dust lattice must be rescaled because discrete site counts inside a sphere exceed its continuum volume at low particle number (Sec.~\ref{sec:dust}). Multi-pass source filters require a sharpening compensation pass \cite{1991ppcs.book.....B} because plain smoothing widens clouds like gravitational softening and makes stable clusters bounce outward (Secs.~\ref{sec:deposition} and \ref{sec:clusters}). Coarser particle clouds---a tempting change on sparse refined meshes---raise constraint violations and drive black hole mass drift, so clouds belong on the finest local level (Sec.~\ref{sec:selfgravity}). A geometric mass measure of the slice also exposed a ghost-zone defect in the host code's finder for apparent horizons, the outermost closed surfaces of zero outgoing light expansion (Secs.~\ref{sec:numerical_methods} and \ref{sec:dust}). With these elements in place, ten thousand macroparticles reproduce in kind dust collapse \cite{Oppenheimer:1939ue}, the sequence of relativistic clusters \cite{1985ApJ...298...34S,1985ApJ...298...58S}, and the migration of an unstable cluster toward the stable branch \cite{Font:2001ew,Wu:2026nrpic} (Secs.~\ref{sec:dust} and \ref{sec:clusters}). Documented differences from Ref.~\cite{Wu:2026nrpic} appear in the collapse times of unstable clusters and in the collapse of one perturbed near-critical model (Sec.~\ref{sec:clusters}).

Section~\ref{sec:system} formulates collisionless matter in a dynamical spacetime and its macroparticle phase-space representation. Section~\ref{sec:method} presents the particle-mesh method: geodesic pushers, metric reconstruction, stress-energy deposition and filtering, spacetime evolution, adaptive mesh refinement, and numerical methods. Section~\ref{sec:testparticles} tests the geodesic pushers, phase mixing, source conservation, and interior particle removal by lapse excision around fixed and evolving black holes. Section~\ref{sec:selfgravity} studies self-gravitating matter at ten thousand macroparticles across spherical dust collapse and the sequence of relativistic clusters. Section~\ref{sec:assembly} describes how the code was assembled and verified, before the final section states our conclusions and the appendix records the scenario parameters.
